\section{Introduction}
Many real-world multi-agent systems are dynamic and continuous, where agents (nodes) interact and exhibit complex behaviors over time. This results in time-evolving node trajectories and dynamic interaction edges. An example is the spread of  COVID-19 in the U.S., where states act as agents and daily migration patterns across states form interaction edges~\cite{huang2021coupled,covidpolicy}. 
Estimating the counterfactual outcomes over time in such %multi-agent dynamical 
systems are crucial for various applications, such as formulating effective policies and designing medical treatment plans~\cite{seedat2022continuous,bica2020crn,bellot2021policy}. 
This can achieve more accurate predictions than non-causal methods by considering the influence of biased confounders. 
Confounders are variables that have influences on treatments and outcomes. 
For example, the health status of the residents in each state (confounders) can impact their level of adherence to the state's policies (treatments), which can influence future confirmed cases/deaths (outcomes). 
%Therefore in observational data, the treatments are not uniform across different confounder values. 
% Observational data shows non-uniform treatments across confounder values~\cite{}. 
Non-causal methods only learn the statistical associations between treatments and outcomes from observational data, which can have non-uniform treatment distributions across confounder values, potentially leading to incorrect predictions such as taking vaccines can increase the number of confirmed cases for each state. Furthermore, causal inference for multi-agent dynamical systems enables effective decision-making by addressing causal questions such as "What if we remove a policy at a specific time" or "What if we change the order of different policies". Therefore, it serves as a promising tool for policymakers. 

Traditionally, the standard approach for causal inference over time is randomized controlled trials (RCTs)~\cite{chalmers1981method}, which can be very costly to obtain and can raise some ethical problems~\cite{seedat2022continuous,bica2020crn}. Thus, researchers have turned to using observational data and employed methods like 
%start to estimate counterfactual outcomes from observational data such as 
linear regression~\cite{robins2000marginal}, recurrent neural networks (RNNs)~\cite{lim2018forecasting,BicaAJS20}, and Transformers~\cite{MelnychukFF22} to estimate counterfactual outcomes with time dependencies. 
However, causal inference for multi-agent dynamical systems presents unique challenges. 

One is that most existing methods~\cite{seedat2022continuous,bica2020crn} assume that nodes are independent, meaning their trajectories are determined solely by their own treatments. 
%Even when considering the influence of neighboring nodes ~\cite{}, it often assumes static interaction edges, which fails to capture the dynamic nature of interactions such as daily population travel patterns between states in the context of COVID-19. 
Some~\cite{anonymous2023cfgode} considers the influence of neighboring nodes but only assumes static interactions among them, which fails to capture situations such as daily population travel patterns between states in the context of COVID-19. 

In casual terms,  influences of neighboring nodes can be categorized into two parts: 1.) time-dependent neighborhood confounding, where a node’s treatment and outcome may be confounded by the covariates of its neighbors.  For example, if cases in neighboring states rise (covariate), a state may implement a vaccine policy (treatment) that affects future confirmed cases/deaths (outcome). 2.) time-dependent interference, where the outcome of a node can be influenced by the treatments of its neighbors. For example, a state may have reduced future cases/deaths (outcome) if neighboring states have implemented a vaccine policy (covariates), as higher vaccination rates within the population flow network give stronger protection.
%are more likely to be reduced. That is, they receive stronger protection through higher vaccination rates among their social networks. 
As the interaction edges evolve along with node trajectories, the challenges lie in predicting the neighbors of each node (edges) and then addressing the time-dependent neighborhood confounding and interference issues.

Another challenge is that current methods lack the ability to capture the continuous and dynamic effects of multiple treatments on such systems. For instance, the impact of a "stay-at-home" policy may be most significant during its initial implementation, and when a "get-vaccine" policy is subsequently introduced, the combined effect of these policies can result in a different outcome. Existing studies often focus on a single treatment~\cite{anonymous2023cfgode,seedat2022continuous} or simply append fixed multi-hot treatment representations when a node receives them. These fixed treatment representations fail to differentiate the influences of the same treatment administered at different times. 


To tackle these challenges, we propose a novel causal inference framework: the \textbf{Ca}usal \textbf{G}raph \textbf{O}rdinary \textbf{D}ifferential \textbf{E}quations (CAG-ODE) to estimate the continuous counterfactual outcome of a multi-agent dynamical system in the presence of multiple treatments and time-varying confounding and interference. Building upon the recent success of graph ordinary differential equations (ODE) in capturing the continuous interaction among agents~\cite{huang2021coupled,luo2023care,LG-ODE,huang2024tango}, our key innovation is to learn time-dependent representations of simultaneous treatments and incorporate them into the ODE function to accurately account for their casual effects on the system. As nodes and edges are jointly evolving, we utilize two coupled treatment-induced ODE functions to account for their respective dynamics. To mitigate confounding bias, we further design two adversarial learning losses, which enable our model to learn balanced continuous trajectory representations unaffected by treatments or interference. Experiments on both real and simulated datasets demonstrate the effectiveness of our proposed model. The primary contributions of this paper can be summarized as follows:
\begin{itemize}
    \item We propose CAG-ODE to estimate continuous counterfactual outcomes in multi-agent systems with evolving interaction edges and multiple treatments.
    \item CAG-ODE features a novel treatment fusing module that can capture the dynamic effects of treatment over time and the combined effect of multiple treatments.
    \item Our method achieves the state-of-art results in counterfactual estimation across varying systems, and can serve as a promising tool for policymakers.
\end{itemize}
